\documentclass[12pt,a4paper]{article}
\usepackage{fontspec}
\setmainfont{Times New Roman}
\usepackage[top=2cm,bottom=2cm,left=2.5cm,right=2cm]{geometry}
\usepackage{enumitem}
\usepackage{tabularx}
\usepackage{booktabs}
\usepackage{array}
\newcolumntype{Y}{>{\centering\arraybackslash}X}
\pagestyle{plain}
\linespread{1.15}

\begin{document}
% --- HEADER 2 CỘT CHUẨN NGHỊ ĐỊNH 30/2020/NĐ-CP ---
\noindent
\begin{minipage}[t]{0.45\textwidth}
    \begin{center}
        \fontsize{11pt}{13pt}\selectfont
        CÔNG AN THÀNH PHỐ HÀ NỘI\\
        \textbf{PHÒNG CẢNH SÁT HÌNH SỰ (PC02)}\\[-0.2em]
        \rule{3.2cm}{0.6pt}\\[0.25cm]
        \fontsize{11pt}{13pt}\selectfont Số: \textbf{07/BB-LK}
    \end{center}
\end{minipage}
\hfill
\begin{minipage}[t]{0.52\textwidth}
    \begin{center}
        \fontsize{11pt}{13pt}\selectfont
        \textbf{CỘNG HÒA XÃ HỘI CHỦ NGHĨA VIỆT NAM}\\
        \textbf{Độc lập - Tự do - Hạnh phúc}\\[-0.2em]
        \rule{3.6cm}{0.6pt}\\[0.25cm]
        \textit{Hà Nội, ngày 25 tháng 07 năm 2016}
    \end{center}
\end{minipage}

\vspace{0.6cm}

\begin{center}
    \textbf{\Large BIÊN BẢN LẤY LỜI KHAI}\\[0.15cm]
    \textit{}
\end{center}

\vspace{0.4cm}
\textit{(Vụ án mạng tại số 14 Đường Bờ Sông, Phân khu Cảng)}



Vào hồi 14 giờ 00 phút, ngày 25 tháng 07 năm 2016, tại Phòng Cảnh sát Hình sự Công an TP. Hà Nội.

\textbf{Chúng tôi gồm:}

\begin{enumerate}[label=\arabic*., leftmargin=1.2cm, itemsep=2pt]
  \item Điều tra viên: Đại úy Lê Minh — Đội Điều tra Trọng án.
  \item Cán bộ ghi biên bản: Trung úy Nguyễn Văn Hoàng — Cán bộ Đội Án mạng.

\end{enumerate}
\textbf{Tiến hành lấy lời khai người liên quan:}

\begin{itemize}[leftmargin=1.2cm, itemsep=2pt]
  \item \textbf{Họ và tên:} \textbf{NGUYỄN NGỌC MAI} | \textbf{Giới tính:} Nữ | \textbf{Sinh ngày:} 15/05/1992 (24 tuổi).
  \item \textbf{CCCD số:} 275190483.
  \item \textbf{Nơi ĐKHKTT:} Số 45, Đường Đoàn Kết, Phường Cảng Đông, TP. Hà Nội.
  \item \textbf{Nghề nghiệp:} Chuyên viên Hành chính — Nhân sự.
  \item \textbf{Mối quan hệ với nạn nhân:} Em họ (con chú ruột của nạn nhân Nguyễn Văn Khang).


\end{itemize}

\subsection*{NỘI DUNG LỜI KHAI (HỎI VÀ ĐÁP)}


\noindent\textbf{Hỏi (ĐTV Lê Minh):} \textit{Chiều và tối ngày 24/07/2016, chị ở đâu, làm gì? Hãy trình bày lý do đến nhà nạn nhân Khang.}\\[0.15cm]
\noindent\textbf{Đáp (Nguyễn Ngọc Mai):} Thưa cán bộ, tôi đến nhà anh Khang để giải quyết tranh chấp quyền sở hữu mảnh đất 200m² tại số 14 Đường Bờ Sông. Khoảng 18:30 chiều 24/07, tôi đi xe máy cùng chồng tôi (Lê Quang Vũ) sang đưa Cam kết tự nguyện trả lại đất ép anh Khang ký nhận. Anh Khang thản nhiên cười thách thức. Trong lúc tức giận, tôi ném toàn bộ tài liệu xuống đất rồi nổ máy xe phóng về nhà trước một mình.\\[0.35cm]

\noindent\textbf{Hỏi (ĐTV Lê Minh):} \textit{Chị có thể trình bày rõ hơn về nguồn gốc mảnh đất 200m² và nguyên nhân dẫn đến mâu thuẫn tranh chấp gay gắt giữa hai gia đình không?}\\[0.15cm]
\noindent\textbf{Đáp (Nguyễn Ngọc Mai):} Mảnh đất 200m² tại số 14 Đường Bờ Sông vốn là ông nội chúng tôi (ông Nguyễn Văn Thọ) để lại cho hai người con trai là bố anh Khang và bố tôi. Khi ông còn sống, anh Khang ở cùng ông. Ba tháng trước sau khi ông nội mất, anh Khang đã lén lút làm Giấy ủy quyền giả mạo chữ ký của ông nội để mang toàn bộ Sổ đỏ mảnh đất 200m² đi thế chấp Ngân hàng vay khoản tiền 1,2 tỷ đồng làm ăn và cho vay nặng lãi. Tuần trước phía Ngân hàng gửi giấy báo nợ về nhà bố tôi ở phố Đoàn Kết thì gia đình tôi mới tá hỏa phát hiện sự việc. Bố tôi và tôi yêu cầu anh Khang phải rút sổ đỏ về, hủy thế chấp và chia trả 50\% diện tích đất (100m²) cho gia đình tôi, nhưng anh Khang chây ì, trơ tráo bảo \textit{"Giấy tờ sang tên đứng chính chủ tao rồi, thích thì cứ kiện lên Tòa"}.\\[0.35cm]

\noindent\textbf{Hỏi (ĐTV Lê Minh):} \textit{Căn cứ vào đâu gia đình chị khẳng định Giấy ủy quyền thế chấp Ngân hàng của nạn nhân Khang là giả mạo?}\\[0.15cm]
\noindent\textbf{Đáp (Nguyễn Ngọc Mai):} Thời điểm hợp đồng thế chấp ký tại Ngân hàng là tháng 04/2016. Lúc đó ông nội tôi đã 86 tuổi, bị tai biến nằm liệt giường tại Bệnh viện từ tháng 02/2016, hoàn toàn mất nhận thức và khả năng hành vi dân sự cho đến khi qua đời. Vậy mà trong hồ sơ vay vốn lại xuất hiện Giấy ủy quyền có xác nhận lăn tay! Rõ ràng anh Khang đã lợi dụng lúc ông mê man để lấy tay ông lăn vào giấy tờ hoặc cấu kết với công chứng viên ngoài để làm giả hồ sơ sang tên Sổ đỏ.\\[0.35cm]

\noindent\textbf{Hỏi (ĐTV Lê Minh):} \textit{Khoản vay thế chấp 1,2 tỷ đồng đó hiện đã đến hạn xử lý ra sao và gây áp lực thế nào lên gia đình chị?}\\[0.15cm]
\noindent\textbf{Đáp (Nguyễn Ngọc Mai):} Phía Ngân hàng đã gửi Thông báo nợ quá hạn lần 3. Họ thông báo nếu trước ngày 30/08/2016 không thanh toán đủ cả gốc lẫn lãi phạt là hơn 1,35 tỷ đồng, Ngân hàng sẽ nộp đơn khởi kiện ra Tòa và tiến hành kê biên, phát mại toàn bộ mảnh đất 200m²! Mảnh đất đó là ông nội để lại, bố tôi dự định lấy 100m² làm nhà thờ họ. Nếu bị Ngân hàng bán đi thì gia đình tôi coi như mất sạch tài sản. Vì vậy bố tôi và tôi vô cùng căm phẫn anh Khang!\\[0.35cm]

\noindent\textbf{Hỏi (ĐTV Lê Minh):} \textit{Trước khi mang đơn sang chiều 24/07, hai bên gia đình đã từng tổ chức thương lượng hay hòa giải lần nào chưa?}\\[0.15cm]
\noindent\textbf{Đáp (Nguyễn Ngọc Mai):} Ngày 15/07/2016, bố tôi cùng bác trưởng họ đã sang nhà số 14 Đường Bờ Sông để hòa giải nội bộ. Gia đình đề nghị anh Khang bán bớt 50m² đất để gom tiền trả nợ Ngân hàng, giải chấp rút Sổ đỏ về rồi chia đều phần đất còn lại. Nhưng anh Khang chỉ trỏ vào mặt bố tôi, thách thức: \textit{"Đất này đứng tên tao trên Sổ đỏ rồi, ông nội cho tao riêng. Mấy người không có quyền gì hết, giỏi thì đâm đơn ra Tòa!"} rồi đuổi bố tôi ra khỏi nhà. Do anh Khang quá ngang ngược và ngoan cố nên tôi mới phải thuê Văn phòng luật sư làm Đơn đòi đất và Cam kết tự nguyện trả lại đất. Nếu anh Khang vẫn nhất quyết không ký giấy trả lại đất thì chúng tôi sẽ khởi kiện.\\[0.35cm]

\noindent\textbf{Hỏi (ĐTV Lê Minh):} \textit{Khi đến nhà Khang lúc 18:30 chiều 24/07, chị đã mang theo những hồ sơ gì và phản ứng của nạn nhân lúc đó ra sao?}\\[0.15cm]
\noindent\textbf{Đáp (Nguyễn Ngọc Mai):} Tôi mang theo 02 bộ hồ sơ gồm: Đơn đòi đất, Cam kết tự nguyện trả lại đất và photocopy Giấy báo nợ của Ngân hàng. Tôi ép anh Khang ký vào bản Cam kết tự nguyện trả lại đất. Nhưng anh Khang cười cợt, bảo đất này ông nội đã cho riêng anh ấy, rồi vứt tài liệu vào mặt tôi. Tôi quá uất ức nên ném tung toàn bộ xấp hồ sơ xuống sàn nhà phòng khách. Anh Vũ bảo tôi cứ đi xe máy về trước, còn anh ấy bận đi nhậu với bạn nên ở lại ra sau.\\[0.35cm]

\noindent\textbf{Hỏi (ĐTV Lê Minh):} \textit{Chị rời nhà nạn nhân lúc mấy giờ? Sau đó chị đi đâu, làm gì?}\\[0.15cm]
\noindent\textbf{Đáp (Nguyễn Ngọc Mai):} Tôi nổ máy xe phóng thẳng một mạch về nhà ở số 45 Phố Đoàn Kết, về đến nhà là khoảng 19:30. Tôi tắm rửa rồi bật tivi lên xem, đang xem thì nhớ tầm khoảng 20h00 - 20h15 tivi bị mất tín hiệu cáp do sự cố đứt cáp ngoài đường. Tôi bực mình tắt tivi vào phòng nằm nghỉ suốt đêm!\\[0.35cm]

\noindent\textbf{Hỏi (ĐTV Lê Minh):} \textit{Chị có quay lại hiện trường hay trực tiếp gây thương tích cho nạn nhân không?}\\[0.15cm]
\noindent\textbf{Đáp (Nguyễn Ngọc Mai):} Tôi thề danh dự là không bao giờ làm chuyện đó! Tôi là người làm hành chính, muốn giải quyết tranh chấp bằng pháp luật nên đã thuê luật sư soạn sẵn hồ sơ định nộp đơn khởi kiện lên Tòa án Nhân dân vào tuần tới.\\[0.35cm]


Biên bản lấy lời khai kết thúc vào hồi 15 giờ 30 phút cùng ngày. Biên bản đã được đọc lại cho người khai nghe, công nhận đúng và cùng ký tên.


\textbf{ĐIỀU TRA VIÊN}

\textit{(Ký tên)}

\textbf{Đại úy Lê Minh}


\textbf{NGƯỜI KHAI BÁO}

\textit{(Ký, ghi rõ họ tên)}

\textbf{Nguyễn Ngọc Mai}

\end{document}
